\scanpage{46}
% Pencil archival foliation: 45.
où disons
\[
\unsure{\gamma_j(X)|\beta_j(X)}\in\mathbb{Z}^{\sigma_j(X)}
\qquad 0\leq j\leq n
\]
Ils sont obtenus en prenant les
$\beta_k(X)$ de Serre [composants de poids purs
de $\operatorname{cl}(X)$] pour $k\geq n$, et \deleted{les composant}
\inserted{en notant}
qu'un tel $\beta_k(X)$ se met sous la forme
\[
\xi^{k-n}\unsure{\beta'_{2n-k}(X)},\quad\text{où}\quad
\unsure{\beta'_{2n-k}(X)}\in\mathrm{M}^{+}_{2n-k}(K);
\]
alors $\gamma_{2n-k}(X)$ s'obtient en prenant la composante
\uline{bipure} de $\unsure{\beta'_{2n-k}(X)}$ ----- .
\margin{\figure{\begin{tikzpicture}
\node[draw,ellipse,align=center] {\unsure{au signe}\\près!};
\end{tikzpicture}}}

Par construction \unsure{$=$|$\overline{\mathrm m}$}, les invariants
ainsi obtenus sont des invariants
$n$-birationnels de $X$. Si $X$ est projective
\uline{lisse sur $K$}\deleted{\unsure{la condition}}, les $\beta'_j(X)$ pour $0\leq j\leq n$
n'est autre (par \uuline{Lefschetz}\ldots) que
la classe des $H^j(X,\unsure{\mathbb{Q}(0)|\mathbb{Z}(0)})$. Donc notre construction
consiste à prendre la composante bipure des
$H^j(X,\unsure{\mathbb{Q}(0)|\mathbb{Z}(0)})$. Mais cela n'est autre aussi.
